\documentclass{amsart}
% For dealing with references we use the comment environment
\usepackage{verbatim}
\newenvironment{reference}{\comment}{\endcomment}
%\newenvironment{reference}{}{}
\newenvironment{slogan}{\comment}{\endcomment}
\newenvironment{history}{\comment}{\endcomment}

% For commutative diagrams we use Xy-pic
\usepackage[all]{xy}

% We use 2cell for 2-commutative diagrams.
\xyoption{2cell}
\UseAllTwocells

% We use multicol for the list of chapters between chapters
\usepackage{multicol}

% This is generally recommended for better output
\usepackage{lmodern}
\usepackage[T1]{fontenc}

% For cross-file-references

% Package for hypertext links:
\usepackage{hyperref}

% For any local file, say "hello.tex" you want to link to please

% Theorem environments.
%
\theoremstyle{plain}
\newtheorem{theorem}[subsection]{Theorem}
\newtheorem{proposition}[subsection]{Proposition}
\newtheorem{lemma}[subsection]{Lemma}

\theoremstyle{definition}
\newtheorem{definition}[subsection]{Definition}
\newtheorem{example}[subsection]{Example}
\newtheorem{exercise}[subsection]{Exercise}
\newtheorem{situation}[subsection]{Situation}

\theoremstyle{remark}
\newtheorem{remark}[subsection]{Remark}
\newtheorem{remarks}[subsection]{Remarks}

\numberwithin{equation}{subsection}

% Macros
%
\def\lim{\mathop{\mathrm{lim}}\nolimits}
\def\colim{\mathop{\mathrm{colim}}\nolimits}
\def\Spec{\mathop{\mathrm{Spec}}}
\def\Hom{\mathop{\mathrm{Hom}}\nolimits}
\def\Ext{\mathop{\mathrm{Ext}}\nolimits}
\def\SheafHom{\mathop{\mathcal{H}\!\mathit{om}}\nolimits}
\def\SheafExt{\mathop{\mathcal{E}\!\mathit{xt}}\nolimits}
\def\Sch{\mathit{Sch}}
\def\Mor{\mathop{\mathrm{Mor}}\nolimits}
\def\Ob{\mathop{\mathrm{Ob}}\nolimits}
\def\Sh{\mathop{\mathit{Sh}}\nolimits}
\def\NL{\mathop{N\!L}\nolimits}
\def\CH{\mathop{\mathrm{CH}}\nolimits}
\def\proetale{{pro\text{-}\acute{e}tale}}
\def\etale{{\acute{e}tale}}
\def\QCoh{\mathit{QCoh}}
\def\Ker{\mathop{\mathrm{Ker}}}
\def\Im{\mathop{\mathrm{Im}}}
\def\Coker{\mathop{\mathrm{Coker}}}
\def\Coim{\mathop{\mathrm{Coim}}}

% Boxtimes
%
\DeclareMathSymbol{\boxtimes}{\mathbin}{AMSa}{"02}

%
% Macros for moduli stacks/spaces
%
\def\QCohstack{\mathcal{QC}\!\mathit{oh}}
\def\Cohstack{\mathcal{C}\!\mathit{oh}}
\def\Spacesstack{\mathcal{S}\!\mathit{paces}}
\def\Quotfunctor{\mathrm{Quot}}
\def\Hilbfunctor{\mathrm{Hilb}}
\def\Curvesstack{\mathcal{C}\!\mathit{urves}}
\def\Polarizedstack{\mathcal{P}\!\mathit{olarized}}
\def\Complexesstack{\mathcal{C}\!\mathit{omplexes}}
% \Pic is the operator that assigns to X its picard group, usage \Pic(X)
% \Picardstack_{X/B} denotes the Picard stack of X over B
% \Picardfunctor_{X/B} denotes the Picard functor of X over B
\def\Pic{\mathop{\mathrm{Pic}}\nolimits}
\def\Picardstack{\mathcal{P}\!\mathit{ic}}
\def\Picardfunctor{\mathrm{Pic}}
\def\Deformationcategory{\mathcal{D}\!\mathit{ef}}

\setlength{\emergencystretch}{3em}
\begin{document}
\title[Stacks Project, Tag 02SX]{459 --- Tensoring a coherent sheaf by a line bundle realizes first-Chern cap in the filtered K quotient}
\maketitle
\noindent Copyright (C) 2005--2025 Johan de Jong. Stacks Project authors. Source: \href{https://stacks.math.columbia.edu/tag/02SX}{Tag 02SX}.

\noindent Editorial difficulty note (not author text). Difficulty H2: Embedded associated points prevent the naive meromorphic-section quotient argument. The proof removes those points through two support-dimension stages, replaces the sheaf by one on its annihilator subscheme, and builds two denominator-ideal quotient sheaves whose cycle difference gives the first-Chern action while remaining in the prescribed K-theory subgroup..

\noindent Editorial provenance note (not author text). History: excerpt from revision \texttt{a0444\allowbreak 6e57e\allowbreak c1fbc\allowbreak 252a8\allowbreak 71afc\allowbreak ec775\allowbreak 2fb28\allowbreak 07b14}, chow.tex, lines 18755--18911. Reference presentation was adapted; original-author context and core support blocks were retained or added. The mathematical wording is unchanged. Licensed under GNU FDL 1.2 or later; no invariant sections or cover texts. See ../licenses/COPYING. Context blocks reproduce author definitions and supporting results; they are not additional counted problems.

\section{Chow groups and K-groups}
\label{section-chow-and-K}

\noindent
In this section we are going to compare $K_0$ of the
category of coherent sheaves to the chow groups.

\medskip\noindent
Let $(S, \delta)$ be as in Situation \ref{situation-setup}.
Let $X$ be a scheme locally of finite type over $S$.
We denote $\textit{Coh}(X) = \textit{Coh}(\mathcal{O}_X)$
the category of coherent sheaves on $X$.
It is an abelian category, see
Cohomology of Schemes, Lemma \href{https://stacks.math.columbia.edu/tag/01Y0}{lemma coherent abelian Noetherian (Tag 01Y0)}.
For any $k \in \mathbf{Z}$ we let $\textit{Coh}_{\leq k}(X)$
be the full subcategory of $\textit{Coh}(X)$
consisting of those coherent sheaves $\mathcal{F}$
having $\dim_\delta(\text{Supp}(\mathcal{F})) \leq k$.



\begin{situation}
\label{situation-setup}
Here $S$ is a locally Noetherian, and universally catenary scheme.
Moreover, we assume $S$ is endowed with a dimension function
$\delta : S \longrightarrow \mathbf{Z}$.
\end{situation}

\begin{lemma}
\label{lemma-cycles-rational-equivalence-K-group}
Let $(S, \delta)$ be as in Situation \ref{situation-setup}.
Let $X$ be a scheme locally of finite type over $S$.
The map
$$
\CH_k(X) \longrightarrow
K_0(\textit{Coh}_{\leq k + 1}(X)/\textit{Coh}_{\leq k - 1}(X))
$$
from Lemma \href{https://stacks.math.columbia.edu/tag/0FDS}{lemma from chow to K (Tag 0FDS)} induces a bijection from
$\CH_k(X)$ onto the image $B_k(X)$ of the map
$$
K_0(\textit{Coh}_{\leq k}(X)/\textit{Coh}_{\leq k - 1}(X))
\longrightarrow
K_0(\textit{Coh}_{\leq k + 1}(X)/\textit{Coh}_{\leq k - 1}(X)).
$$
\end{lemma}

\begin{lemma}
\label{lemma-coherent-sheaf-cap-c1}
Let $(S, \delta)$ be as in Situation \ref{situation-setup}.
Let $X$ be locally of finite type over $S$.
Let $\mathcal{L}$ be an invertible $\mathcal{O}_X$-module.
Let $\mathcal{F}$ be a coherent $\mathcal{O}_X$-module.
Assume $\dim_\delta(\text{Supp}(\mathcal{F})) \leq k + 1$.
Then the element
$$
[\mathcal{F} \otimes_{\mathcal{O}_X} \mathcal{L}]
-
[\mathcal{F}]
\in
K_0(\textit{Coh}_{\leq k + 1}(X)/\textit{Coh}_{\leq k - 1}(X))
$$
lies in the subgroup $B_k(X)$ of
Lemma \ref{lemma-cycles-rational-equivalence-K-group} and maps to
the element $c_1(\mathcal{L}) \cap [\mathcal{F}]_{k + 1}$ via
the map $B_k(X) \to \CH_k(X)$.
\end{lemma}

\begin{proof}
Let
$$
0 \to \mathcal{K} \to \mathcal{F} \to \mathcal{F}' \to 0
$$
be the short exact sequence constructed in
Divisors, Lemma \href{https://stacks.math.columbia.edu/tag/02OL}{lemma remove embedded points (Tag 02OL)}.
This in particular means that $\mathcal{F}'$ has no embedded
associated points.
Since the support of $\mathcal{K}$ is nowhere dense in the
support of $\mathcal{F}$ we see that
$\dim_\delta(\text{Supp}(\mathcal{K})) \leq k$. We may
re-apply
Divisors, Lemma \href{https://stacks.math.columbia.edu/tag/02OL}{lemma remove embedded points (Tag 02OL)}
starting with $\mathcal{K}$ to get a short exact sequence
$$
0 \to \mathcal{K}'' \to \mathcal{K} \to \mathcal{K}' \to 0
$$
where now $\dim_\delta(\text{Supp}(\mathcal{K}'')) < k$
and $\mathcal{K}'$ has no embedded associated points.
Suppose we can prove the lemma for the coherent sheaves
$\mathcal{F}'$ and $\mathcal{K}'$. Then we see
from the equations
$$
[\mathcal{F}]_{k + 1}
=
[\mathcal{F}']_{k + 1}
+ [\mathcal{K}']_{k + 1}
+ [\mathcal{K}'']_{k + 1}
$$
(use Lemma \href{https://stacks.math.columbia.edu/tag/02QZ}{lemma additivity sheaf cycle (Tag 02QZ)}),
$$
[\mathcal{F} \otimes_{\mathcal{O}_X} \mathcal{L}]
-
[\mathcal{F}]
=
[\mathcal{F}' \otimes_{\mathcal{O}_X} \mathcal{L}]
-
[\mathcal{F}']
+
[\mathcal{K}' \otimes_{\mathcal{O}_X} \mathcal{L}]
-
[\mathcal{K}']
+
[\mathcal{K}'' \otimes_{\mathcal{O}_X} \mathcal{L}]
-
[\mathcal{K}'']
$$
(use the $\otimes \mathcal{L}$ is exact)
and the trivial vanishing of $[\mathcal{K}'']_{k + 1}$ and
$[\mathcal{K}'' \otimes_{\mathcal{O}_X} \mathcal{L}]
- [\mathcal{K}'']$ in
$K_0(\textit{Coh}_{\leq k + 1}(X)/\textit{Coh}_{\leq k - 1}(X))$
that the result holds
for $\mathcal{F}$. What this means is that we may assume that
the sheaf $\mathcal{F}$ has no embedded associated points.

\medskip\noindent
Assume $X$, $\mathcal{F}$ as in the lemma, and assume in addition
that $\mathcal{F}$ has no embedded associated points. Consider the
sheaf of ideals $\mathcal{I} \subset \mathcal{O}_X$, the corresponding
closed subscheme $i : Z \to X$ and the coherent $\mathcal{O}_Z$-module
$\mathcal{G}$ constructed in
Divisors, Lemma \href{https://stacks.math.columbia.edu/tag/02OM}{lemma no embedded points endos (Tag 02OM)}.
Recall that $Z$ is a locally Noetherian scheme without embedded points,
$\mathcal{G}$ is a coherent sheaf without embedded
associated points, with $\text{Supp}(\mathcal{G}) = Z$
and such that $i_*\mathcal{G} = \mathcal{F}$.
Moreover, set $\mathcal{N} = \mathcal{L}|_Z$.

\medskip\noindent
By Divisors, Lemma \href{https://stacks.math.columbia.edu/tag/02OZ}{lemma regular meromorphic section exists (Tag 02OZ)}
the invertible sheaf $\mathcal{N}$ has a regular meromorphic section $s$
over $Z$. Let us denote $\mathcal{J} \subset \mathcal{O}_Z$ the sheaf
of denominators of $s$. By Lemma \href{https://stacks.math.columbia.edu/tag/02SW}{lemma no embedded points (Tag 02SW)}
there exist short exact sequences
$$
\begin{matrix}
0 &
\to &
\mathcal{J}\mathcal{G} &
\xrightarrow{1} &
\mathcal{G} &
\to &
\mathcal{Q}_1 &
\to &
0 \\
0 &
\to &
\mathcal{J}\mathcal{G} &
\xrightarrow{s} &
\mathcal{G} \otimes_{\mathcal{O}_Z} \mathcal{N} &
\to &
\mathcal{Q}_2 &
\to &
0
\end{matrix}
$$
such that $\dim_\delta(\text{Supp}(\mathcal{Q}_i)) \leq k$ and
such that the cycle
$
[\mathcal{Q}_2]_k - [\mathcal{Q}_1]_k
$
is a representative of $c_1(\mathcal{N}) \cap [\mathcal{G}]_{k + 1}$.
We see (using the fact that
$i_*(\mathcal{G} \otimes \mathcal{N}) = \mathcal{F} \otimes \mathcal{L}$
by the projection formula, see
Cohomology, Lemma \href{https://stacks.math.columbia.edu/tag/01E8}{lemma projection formula (Tag 01E8)})
that
$$
[\mathcal{F} \otimes_{\mathcal{O}_X} \mathcal{L}]
-
[\mathcal{F}]
=
[i_*\mathcal{Q}_2] - [i_*\mathcal{Q}_1]
$$
in $K_0(\textit{Coh}_{\leq k + 1}(X)/\textit{Coh}_{\leq k - 1}(X))$.
This already shows that
$[\mathcal{F} \otimes_{\mathcal{O}_X} \mathcal{L}] - [\mathcal{F}]$
is an element of $B_k(X)$. Moreover we have
\begin{eqnarray*}
[i_*\mathcal{Q}_2]_k - [i_*\mathcal{Q}_1]_k
& = &
i_*\left( [\mathcal{Q}_2]_k - [\mathcal{Q}_1]_k \right) \\
& = &
i_*\left(c_1(\mathcal{N}) \cap [\mathcal{G}]_{k + 1} \right) \\
& = &
c_1(\mathcal{L}) \cap i_*[\mathcal{G}]_{k + 1} \\
& = &
c_1(\mathcal{L}) \cap [\mathcal{F}]_{k + 1}
\end{eqnarray*}
by the above and Lemmas \href{https://stacks.math.columbia.edu/tag/02SU}{lemma pushforward cap c1 (Tag 02SU)}
and \href{https://stacks.math.columbia.edu/tag/02R6}{lemma cycle push sheaf (Tag 02R6)}. And this agree with the
image of the element under $B_k(X) \to \CH_k(X)$ by definition.
Hence the lemma is proved.
\end{proof}
\end{document}
